\section{Six lettered conditions share one reader and one prompt}
\label{sec:setup}
Every condition answers the same question at step $t$: given the task and the current screen's OCR text, what text should the agent type? Memory conditions run in two stages (B merges them into one call). The reader first \emph{localizes}, naming one earlier screen that it believes holds the needed text, then \emph{reads} the full OCR text of that screen and produces the text to type. The reader is Qwen2.5-14B-Instruct~\citep{qwen25}, greedy, with one shared prompt across conditions; \S\ref{sec:length} repeats the main comparison with the 7B reader.

\paragraph{Conditions.}
\textbf{A, no history} sees only the task and the current screen. \textbf{B, full history} sees the OCR text of every earlier screen (1,241 tokens on average; 4,584 for the longest GUI-Odyssey trajectory) and names screen and text in one call. \textbf{C, summary} replaces each screen's text with one LLM-written summary (250 tokens in total). \textbf{D, directory-only} sees one index line per screen (353 tokens) and answers from the lines alone. \textbf{E, screen index} sees the same lines, names a screen, then reads it in full; D and E share the localization stage. E is ReadAgent's gist-and-lookup~\citep{readagent} on screens, the method we set out to improve. \textbf{F, gold-screen read} reads the gold source screen and bounds the reading stage. \textbf{BM25} localizes by lexical retrieval with the task plus current screen as query, or with the gold text as query (\textbf{BM25, gold query}). \textbf{Gating rule}, our own addition, reads only the two previous screens when their text overlaps the current one and otherwise falls back to E; we built it on the unfiltered long set and report it only in \S\ref{sec:pitfalls}.

\paragraph{Writer.}
Qwen2.5-7B-Instruct writes each index line from the screen's OCR text in at most 32 tokens, naming the screen and the kinds of field it shows without copying values (round \textbf{r0}). One training round and four prompt variants each change one thing. \textbf{r1} applies direct preference optimization~\citep{dpo} to 224 pairs in which the preferred line led the reader to a valid screen and the rejected line did not. The pairs come from the oracle-index candidates of 47 source screens of the evaluation cases themselves, a choice that favors r1. \textbf{r2} forces a two-field format, \textbf{r3} permits values in the line, \textbf{mv3} (multi-view) writes three lines per screen from three prompts, and \textbf{rp} rewrites every screen's line with the oracle's retrieval prompt and sampling. These four keep the r0 weights. The \textbf{oracle index} replaces only the source screen's line with one of five candidates: the r0 line and four sampled at temperature 0.9 from a retrieval prompt. It counts a case as localized if any candidate leads the reader to a valid screen.

\paragraph{Metrics and tests.}
\emph{Localization} is correct when the named screen is any screen whose OCR text contains the gold text (\emph{any-valid scoring}; \S\ref{sec:pitfalls} quantifies the alternative). \emph{End-to-end} is correct when the produced text, normalized as in mining, contains the gold text. We pair all comparisons on the same cases and test them with McNemar's test, reporting the net number of cases that flip and the two-sided $p$.
